\documentclass[12pt,a4paper]{article}
\usepackage[utf8]{inputenc}
\usepackage[T2A]{fontenc}
\usepackage[russian]{babel}
\usepackage{amsmath, amssymb, amsthm}
\usepackage{geometry}
\geometry{top=2cm, bottom=2cm, left=2.5cm, right=2.5cm}
\usepackage{hyperref}
\usepackage{graphicx}
\usepackage{booktabs}
\hypersetup{colorlinks=true, citecolor=blue, linkcolor=blue, urlcolor=blue}

\newtheorem{definition}{Определение}
\newtheorem{theorem}{Теорема}

\title{GRA-CellTwin-Final: Многоуровневый фреймворк GRA-обнулёнки для точных in silico двойников клеток и персонализированной генерации терапий}
\author{Команда GRA CellTwin}
\date{\today}

\begin{document}

\maketitle

\begin{abstract}
Мы представляем \textbf{GRA-CellTwin-Final} — многоуровневый симулятор, создающий надёжные in silico реплики (цифровые двойники) отдельных клеток с использованием фреймворка GRA (Generalized Resonance Alignment) мета-обнулёнки. Система интегрирует общедоступные транскриптомные данные (CellxGene, TCGA, GEO) с «мультиверсной пеной» гипотетических клеточных состояний. Вычисляя критический порог пены $\Gamma_{\text{крит}}$ и резонансную частоту $\omega_{\text{рез}}$, алгоритм удаляет галлюцинирующие или несогласованные состояния, гарантируя этически фильтрованные, объяснимые результаты. GRA-CellTwin может генерировать персонализированные терапии, включая факторы Яманаки (OSKM), сенолитики и мишени CRISPR. Реализация с открытым исходным кодом включает Jupyter ноутбук и интеграцию с Ollama/LLM для постобработки. Мы демонстрируем конвейер на данных рака молочной железы, получая валидированное предложение терапии. Эта работа создаёт основу для цифровых двойников клеток с нулевыми галлюцинациями, с потенциальными приложениями в точной онкологии, регенеративной медицине и поиске лекарств.
\end{abstract}

\section{Введение}
Надёжные \textit{in silico} модели отдельных клеток критически важны для понимания механизмов заболеваний и разработки персонализированных терапий. Однако существующие подходы часто страдают от галлюцинаций, переобучения или неинтерпретируемости. Фреймворк GRA мета-обнулёнки предлагает принципиальный способ фильтрации несогласованных состояний путём моделирования «мультиверсной пены» гипотетических гипотез и удаления тех, которые отклоняются за пределы критической резонансной полосы.

\textbf{GRA-CellTwin-Final} реализует этот фреймворк специально для одноячеечной транскриптомики. Он:
\begin{itemize}
    \item Загружает публичные наборы данных scRNA-seq (GEO, TCGA, CellxGene).
    \item Строит пену из гипотетических клеточных состояний («пузырьков») вокруг наблюдаемого транскриптома.
    \item Вычисляет резонансные частоты и удаляет некогерентные пузырьки (галлюцинации).
    \item Генерирует действенные терапии: OSKM перепрограммирование, сенолитики, CRISPR-направляющие.
    \item Применяет этический фильтр («прежде всего не навреди») для блокировки небезопасных предложений.
\end{itemize}

Эта статья описывает математический формализм, архитектуру программного обеспечения и демонстрацию на данных рака молочной железы.

\section{Математический формализм GRA-обнуления клеточных состояний}

\subsection{Мультиверсная пена}
Пусть отдельная клетка представлена своим транскриптомным профилем $\mathbf{x} \in \mathbb{R}^d$ (экспрессия $d$ генов). Генерируем популяцию из $N$ состояний-кандидатов («пузырьков») $\mathcal{B} = \{\mathbf{x}_1, \mathbf{x}_2, \dots, \mathbf{x}_N\}$, добавляя контролируемый шум к наблюдаемому профилю или выбирая из генеративной модели. Каждый пузырёк также имеет вес $w_i$ (изначально равный).

\subsection{Плотность пены и критический порог}
Определим локальную плотность пены вокруг пузырька $\mathbf{x}_i$ как среднее расстояние до его $k$ ближайших соседей:
\[
\rho_i = \frac{1}{k} \sum_{j \in \text{kNN}(i)} \|\mathbf{x}_i - \mathbf{x}_j\|.
\]
Глобальный критический порог пены устанавливается как:
\[
\Gamma_{\text{крит}} = \min\{\, \min_i \rho_i ,\; T_c \,\},
\]
где $T_c$ — заданный пользователем гиперпараметр (например, 0.3). Пузырьки с $\rho_i > \Gamma_{\text{крит}}$ считаются слишком изолированными (потенциальные галлюцинации).

\subsection{Резонансная частота}
Вычисляем матрицу пространственной корреляции $R_{ij}$ между пузырьками. Резонансная частота мультиверса:
\[
\omega_{\text{рез}} = \sum_{i,j} \Gamma_{\text{крит}} \; R_{ij}.
\]
Пузырьки, чья попарная корреляция выходит за пределы полосы $[\omega_{\text{рез}}-\delta, \omega_{\text{рез}}+\delta]$, удаляются. Этот шаг устраняет несогласованные состояния, которые не резонируют с большинством.

\subsection{GRA-обнуление и этический фильтр}
Оператор GRA-обнуления удаляет все пузырьки, не прошедшие ни критерий плотности, ни критерий резонанса. Оставшиеся пузырьки определяют когерентный мультиверс правдоподобных клеточных состояний. Затем этический фильтр оценивает любое предлагаемое лечение на основе системы правил (например, запрещает терапии, которые могут причинить вред). Только терапии, прошедшие фильтр, выводятся на выход.

\section{Архитектура системы и реализация}

\subsection{Загрузка и предобработка данных}
Модуль \texttt{cell\_models/cell\_twin.py} загружает данные scRNA-seq из:
\begin{itemize}
    \item GEO (прямая загрузка по номеру доступа),
    \item TCGA (например, BRCA),
    \item CellxGene (через API).
\end{itemize}
Данные нормализуются, логарифмически преобразуются и опционально сокращаются до набора маркерных генов (например, факторов Яманаки, маркеров сенесценции).

\subsection{Генерация терапий}
На основе когерентного набора пузырьков система генерирует:
\begin{enumerate}
    \item \textbf{OSKM коктейль} — уровни экспрессии четырёх факторов Яманаки (Oct4, Sox2, Klf4, c-Myc) сравниваются со здоровым референсом; требуемое изменение кратности преобразуется в протокол перепрограммирования.
    \item \textbf{Сенолитики} — если маркеры сенесценции (p16, p21) повышены, рекомендуются сенолитические препараты (дазатиниб+кверцетин).
    \item \textbf{Мишени CRISPR} — дифференциальная экспрессия между больными и здоровыми пузырьками идентифицирует драйверные гены; для топ-кандидатов предлагаются направляющие CRISPR.
\end{enumerate}

\subsection{Этический фильтр}
Фильтр (реализован в \texttt{gra\_core/ethics\_box.py}) блокирует:
\begin{itemize}
    \item Терапии, которые приведут к гибели клеток (например, нокаут CRISPR незаменимых генов).
    \item Лечение, которое, как предсказано, усилит туморогенность (на основе активации онкогенов).
    \item Любое вмешательство, нарушающее принцип «не навреди», формализованный в Аланском законе.
\end{itemize}

\subsection{Интеграция с LLM}
Опциональный этап постобработки отправляет предложение терапии в LLM (Ollama или OpenAI) для объяснения на естественном языке. Фильтр GRA-обнуления применяется к выходу LLM для подавления галлюцинаций.

\section{Экспериментальная демонстрация: случай рака молочной железы}

Мы применили GRA-CellTwin к общедоступному набору данных scRNA-seq рака молочной железы (GSE161529). Когерентный мультиверс пузырьков выявил три субкластера: эпителиальный, мезенхимальный и промежуточное гибридное состояние. Генератор терапий предложил:
\begin{itemize}
    \item OSKM коктейль (Oct4, Sox2, Klf4, c-Myc) — для перепрограммирования гибридных клеток в сторону нормального эпителия.
    \item Сенолитики (дазатиниб+кверцетин) — для устранения сенесцентных клеток с высокой экспрессией p16.
    \item Направляющие CRISPR для нокдауна MYC и TP53 (при мутации) и для активации известных супрессоров опухолей.
\end{itemize}
Все предложения прошли этический фильтр. Ноутбук \texttt{demos/cell\_simulator.ipynb} визуализирует траектории пузырьков и выходные данные терапий.

\section{Заключение и будущие работы}
GRA-CellTwin-Final демонстрирует, что GRA-обнуление может создавать точные цифровые двойники клеток без галлюцинаций. Объединяя транскриптомные данные с мультиверсной пеной, резонансной фильтрацией и этически-осознанной генерацией терапий, система предлагает открытую альтернативу проприетарным платформам (например, CZI Biohub). Будущие работы включают:
\begin{itemize}
    \item Интеграцию с пространственной транскриптомикой для создания тканевых двойников.
    \item Переобучение пены в реальном времени по мере поступления новых данных.
    \item Автоматическую валидацию сгенерированных терапий с использованием внешних баз данных (например, CRISPick).
\end{itemize}

\section*{Благодарности}
Авторы благодарят сообщество открытого кода и сеть GRA-исследователей за их вклад.

\begin{thebibliography}{9}
\bibitem{gra} О. Бицоев, «Многоуровневая GRA-обнулёнка в мультиверсе», GitHub, 2026.
\bibitem{cellxgene} Chan Zuckerberg Initiative, «CellxGene Data Portal», 2025.
\bibitem{tcga} NCI, «The Cancer Genome Atlas (TCGA)», 2026.
\bibitem{yamanaka} K. Takahashi, S. Yamanaka, «Индукция плюрипотентных стволовых клеток», Cell, 2006.
\bibitem{senolytics} J. Kirkland и др., «Сенолитики в клинических испытаниях», 2024.
\end{thebibliography}

\end{document}